\documentclass[12pt]{beamer}

\usepackage[brazil]{babel}
\usepackage[T1]{fontenc}
\usepackage[utf8]{inputenc}
\usepackage{enumitem, tikz}

\setitemize{
  itemsep=1em,
  label=\usebeamerfont*{itemize item}
    \usebeamercolor[fg]{itemize item}
    \usebeamertemplate{itemize item}
}

\usetikzlibrary{positioning, decorations.markings}
\tikzset{small/.style={draw,fill,circle,inner sep=1pt,outer sep=0pt}}

\setbeamertemplate{footline}[frame number]{}
\setbeamertemplate{navigation symbols}{}

\title{Introdução à teoria da ordem}
\author{Gustavo Zambonin}
\institute{
  \includegraphics[scale=0.15]{ufsc}                    \\ \vspace{-4mm}
  Universidade Federal de Santa Catarina                \\
  Departamento de Informática e Estatística             \\
  INE5601 --- Fundamentos Matemáticos da Informática    \\ \vspace{2mm}
  \texttt{gustavo.zambonin@posgrad.ufsc.br}
}
\date{}

\begin{document}

\begin{frame}[plain,noframenumbering]
  \titlepage
\end{frame}

\begin{frame}
  \frametitle{Contexto}
  \begin{itemize}
    \item<1-> Relação: subconjunto do produto Cartesiano de dois (ou mais)
        conjuntos quaisquer
    \begin{itemize}
      \item $\equiv_n$, a congruência módulo $n$ em $\mathbb{Z}$, ou seja,
          $x \equiv y \pmod{n}$
    \end{itemize}
    \item<2-> Propriedades de relações: reflexividade, irreflexividade,
        simetria, assimetria, antissimetria, transitividade
    \begin{itemize}
      \item $\equiv_n$ é uma relação de equivalência
    \end{itemize}
    \item<3-> É intuitivo pensar que relações são utilizadas para construir
        ordenamentos
  \end{itemize}
\end{frame}

\begin{frame}
  \frametitle{Exemplos práticos}
  \begin{itemize}
    \item Organização de verbetes em um dicionário
    \begin{itemize}
      \item Letras no alfabeto e tamanho do verbete
    \end{itemize}
    \item Descrição de um grafo curricular
    \begin{itemize}
      \item Matérias como pré-requisitos
    \end{itemize}
    \item Parentesco entre pessoas, escalonamento temporal de tarefas etc.
    \item<2-> Ideia geral: comparação de (alguns) elementos entre conjuntos
        quaisquer
  \end{itemize}
\end{frame}

\begin{frame}
  \frametitle{Propriedades de relações}
  \begin{itemize}
    \item Considere um conjunto $S$ e uma relação $R$, construída a partir de
        $S \times S$
    \item $R$ é considerada uma relação \emph{parcial} se é reflexiva,
        antissimétrica e transitiva
    \item<1-> Ou seja, $\forall a, b, c \in S$,
    \begin{itemize}[itemsep=0pt]
      \item<2-> $(a, a) \in R$ (reflexividade)
      \item<3-> $(a, b), (b, a) \in R \Rightarrow a = b$ (antissimetria)
      \item<4-> $(a, b), (b, c) \in R \Rightarrow (a, c) \in R$ (transitividade)
    \end{itemize}
  \end{itemize}
\end{frame}

\begin{frame}
  \frametitle{Conjuntos parcialmente ordenados}
  \begin{itemize}
    \item Quando parcial, $R$ é geralmente denotada com o símbolo da relação,
        ou $\leq$, $\preccurlyeq$
    \item O par ordenado $(S, \preccurlyeq)$ é então chamado de conjunto
        parcialmente ordenado, ou \textbf{poset} (\emph{partially ordered set})
    \item $a, b$ são \textbf{comparáveis} se $a \preccurlyeq b$ ou $b
        \preccurlyeq a$, e incomparáveis caso contrário
    \begin{itemize}
      \item se $a \neq b$, a notação $a \prec b$ é utilizada
    \end{itemize}
  \end{itemize}
\end{frame}

\begin{frame}
  \frametitle{Exemplo}
  \framesubtitle{Conjuntos parcialmente ordenados}
  \begin{itemize}
    \item<1-> Considere o conjunto dos inteiros $\mathbb{Z}$ e a relação de
        ``menor ou igual'' $\leq$. $(\mathbb{Z}, \leq)$ é um \emph{poset}?
    \item<2-> Tome inteiros quaisquer $a, b, c \in \mathbb{Z}$
    \begin{itemize}[itemsep=0pt]
      \item<3-> $a \leq a$ para todo inteiro, portanto $\leq$ é reflexiva
      \item<4-> $a \leq b, b \leq a \Rightarrow a = b$, portanto $\leq$ é
          antissimétrica
      \item<5-> $a \leq b, b \leq c \Rightarrow a \leq c$, portanto $\leq$ é
          transitiva
    \end{itemize}
    \item<6-> Logo, $(\mathbb{Z}, \leq)$ é um \emph{poset}
  \end{itemize}
\end{frame}

\begin{frame}
  \frametitle{Exemplo}
  \framesubtitle{Conjuntos parcialmente ordenados}
  \begin{itemize}
    \item<1-> Considere um conjunto $T$ e seu conjunto potência,
        $\mathcal{P}(T)$, e a relação de ``subconjunto'' $\subseteq$.
          $(\mathcal{P}(T), \subseteq)$ é um \emph{poset}?
    \item<2-> Tome conjuntos quaisquer $t_1, t_2, t_3 \in \mathcal{P}(T)$
    \begin{itemize}[itemsep=0pt]
      \item<3-> $t_1 \subseteq t_1$ para todo conjunto, portanto $\subseteq$ é
          reflexiva
      \item<4-> $t_1 \subseteq t_2, t_2 \subseteq t_1 \Rightarrow t_1 = t_2$,
          portanto $\subseteq$ é antissimétrica
      \item<5-> $t_1 \subseteq t_2, t_2 \subseteq t_3 \Rightarrow t_1 \subseteq
          t_3$, portanto $\subseteq$ é transitiva
    \end{itemize}
    \item<6-> Logo, $(\mathcal{P}(T), \subseteq)$ é um \emph{poset}
  \end{itemize}
\end{frame}

\begin{frame}
  \frametitle{Exemplo}
  \framesubtitle{Conjuntos parcialmente ordenados}
  \begin{itemize}
    \item<1-> Considere o conjunto de seres humanos $H$, e a relação de ``é
        mais velho que'' $\preccurlyeq$. $(H, \preccurlyeq)$ é um \emph{poset}?
    \item<2-> Tome pessoas quaisquer $h_1, h_2, h_3 \in H$
    \begin{itemize}[itemsep=0pt]
      \item<3-> $h_1 \preccurlyeq h_2 \Rightarrow h_2 \not\preccurlyeq h_1$
          (antissimetria: se $h_1$ é mais velha que $h_2$, o oposto não pode
            ser verdade)
      \item<4-> $h_1 \preccurlyeq h_2, h_2 \preccurlyeq h_3 \Rightarrow h_1
          \preccurlyeq h_3$ (transitividade: se $h_1$ é mais velha que $h_2$, e
            $h_2$ é mais velha que $h_3$, então $h_1$ é mais velha que $h_3$)
      \item<5-> $h_1 \not\preccurlyeq h_1$ (reflexividade: uma pessoa não pode
          ser mais velha do que si mesma)
    \end{itemize}
    \item<6-> Portanto, $(H, \preccurlyeq)$ não é um \emph{poset}
  \end{itemize}
\end{frame}

\begin{frame}
  \frametitle{Conjuntos parcialmente ordenados estritos}
  \begin{itemize}
    \item Em alguns contextos, a definição de \emph{poset} apresentada é
        chamada de ``conjunto parcialmente ordenado não-estrito''
    \item<1-> O par ordenado $(S, R)$ é um conjunto parcialmente ordenado
        estrito se, $\forall a, b, c \in S$,
    \begin{itemize}[itemsep=0pt]
      \item<2-> $(a, a) \not\in R$ (irreflexividade)
      \item<3-> $(a, b), (b, c) \in R \Rightarrow (a, c) \in R$ (transitividade)
    \end{itemize}
    \item Para transformar um \emph{poset} em um \emph{poset} estrito,
        remova todos os pares $(a, a) \in R$
    \item A notação para relações parciais estritas arbitrárias é similar,
        utilizando os símbolos $<, \prec$
  \end{itemize}
\end{frame}

\begin{frame}
  \frametitle{Exemplo}
  \framesubtitle{Conjuntos parcialmente ordenados estritos}
  \begin{itemize}
    \item<1-> Considere o conjunto de matérias em um currículo $C$, e a relação
        de ``é pré-requisito de'' $\prec$. $(C, \prec)$ é um \emph{poset}
          estrito?
    \item<2-> Tome matérias quaisquer $c_1, c_2, c_3 \in C$
    \begin{itemize}[itemsep=0pt]
      \item<3-> $c_1 \not\prec c_1$ (irreflexividade: $c_1$ não pode ser
          pré-requisito dela mesma)
      \item<4-> $c_1 \prec c_2, c_2 \prec c_3 \Rightarrow c_1 \prec c_3$
          (transitividade: se $c_1$ é pré-requisito de $c_2$, e $c_2$ é
            pré-requisito de $c_3$, então $c_1$ é pré-requisito de $c_3$)
    \end{itemize}
    \item<5-> Portanto, $(C, \prec)$ é um \emph{poset} estrito
  \end{itemize}
\end{frame}

\begin{frame}
  \frametitle{Dualidade de \emph{posets}}
  \begin{itemize}
    \item Para qualquer relação binária $R$, uma relação inversa $R^{-1}$ pode
        ser construída trocando a ordem dos elementos de todos os pares
          ordenados em $R$
    \begin{itemize}
      \item As propriedades usuais de relações mantêm-se em $R^{-1}$
    \end{itemize}
    \item Formalmente, para dois conjuntos quaisquer $S, T$, $R^{-1} = \{(t, s)
        \in T \times S \mid (s, t) \in R\}$
    \item No caso de relações parciais, os símbolos $\succcurlyeq, \geq, \succ,
        >$ podem ser utilizados para denotar a inversa, neste caso chamada de
          \textbf{dual}
  \end{itemize}
\end{frame}

\begin{frame}
  \frametitle{Exemplo}
  \framesubtitle{Conjuntos parcialmente ordenados duais}
  \begin{itemize}
    \item<1-> Considere o conjunto dos inteiros $\mathbb{Z}$ e a relação de
        ``maior ou igual'' $\geq$
    \begin{itemize}
      \item<2-> $(\mathbb{Z}, \geq)$ é um \emph{poset}
    \end{itemize}
    \item<3-> Considere um conjunto $T$ e seu conjunto potência,
        $\mathcal{P}(T)$, e a relação de ``superconjunto'' $\supseteq$
    \begin{itemize}
      \item<4-> $(\mathcal{P}(T), \supseteq)$ é um \emph{poset}
    \end{itemize}
    \item<5-> Considere o conjunto de matérias em um currículo $C$, e a relação
        de ``é pós-requisito de'' $\succ$
    \begin{itemize}
      \item<6-> $(C, \succ)$ é um \emph{poset} estrito
    \end{itemize}
  \end{itemize}
\end{frame}

\begin{frame}
  \frametitle{Conjuntos totalmente ordenados}
  \begin{itemize}
    \item Quando todos os elementos de uma relação parcial são comparáveis,
        esta é chamada de \textbf{ordem total}
    \item<1-> O par ordenado $(S, R)$ é chamado de conjunto totalmente
        ordenado, ou \textbf{toset} (\emph{totally ordered set}) se, $\forall
          a, b \in S$,
    \begin{itemize}[itemsep=0pt]
      \item<2-> $R$ é uma relação parcial
      \item<3-> Ou $(a, b) \in R$, ou $(b, a) \in R$, ou $a = b$ (tricotomia)
    \end{itemize}
    \item São um caso particular de \emph{posets}, e portanto, a dualidade
        mantém-se
    \item Também podem ser chamados de \textbf{cadeias} e ordens lineares, e
        seus duais, de anticadeias
  \end{itemize}
\end{frame}

\begin{frame}
  \frametitle{Exemplo}
  \framesubtitle{Conjuntos totalmente ordenados}
  \begin{itemize}
    \item<1-> Considere o conjunto dos inteiros positivos $\mathbb{Z}^{+}$ e a
        relação de divisibilidade $\mid$. $(\mathbb{Z}^{+}, \mid)$ é um
          \emph{poset}? Se sim, é também um \emph{toset}?
    \item<2-> Para quaisquer inteiros $a, b, c \in \mathbb{Z}^{+}$, e quaisquer
        dois primos $p_1, p_2 \in \mathbb{Z}^{+}$
    \begin{itemize}[itemsep=0pt]
      \item<3-> $a \mid a$ (reflexividade)
      \item<4-> $a \mid b, b \mid a \Rightarrow a = b$ (antissimetria)
      \item<5-> $a \mid b, b \mid c \Rightarrow a \mid c$ (transitividade)
      \item<6-> $p_1 \nmid p_2, p_2 \nmid p_1$ (existem elementos incomparáveis)
    \end{itemize}
    \item<7-> Então, $(\mathbb{Z}^{+}, \mid)$ é um \emph{poset}, mas não um
        \emph{toset}
  \end{itemize}
\end{frame}

\begin{frame}
  \frametitle{Exemplo}
  \framesubtitle{Conjuntos totalmente ordenados}
  \begin{itemize}
    \item<1-> O \emph{poset} $(\mathbb{Z}, \leq)$ é também um \emph{toset},
        pois quaisquer inteiros podem ser comparados com $\leq$
    \item<2-> O \emph{poset} $(\mathcal{P}(T), \subseteq)$ não é um \emph{toset}
    \begin{itemize}
      \item Considere $T = \{x, y\}, \{x\}, \{y\} \in \mathcal{P}(T)$ e veja
          que $\{x\} \not\subseteq \{y\}, \{y\} \not\subseteq \{x\}$
    \end{itemize}
    \item<3-> O produto Cartesiano de qualquer número de \emph{tosets},
        utilizando ordenação lexicográfica, é um \emph{toset}
    \begin{itemize}[itemsep=0pt]
      \item Dados dois \emph{tosets} $(P, \preccurlyeq_1), (Q,
          \preccurlyeq_2)$, e elementos quaisquer $p_1, p_2 \in P, q_1, q_2 \in
            Q$
      \item Para $(P_1 \times P_2, \preccurlyeq)$, a relação $(p_1, q_1)
          \preccurlyeq (p_2, q_2) \Rightarrow p_1 \prec_1 p_2 \text{ ou } (p_1
            = p_2 \text{ e } q_1 \preccurlyeq_2 q_2)$
    \end{itemize}
  \end{itemize}
\end{frame}

\begin{frame}
  \frametitle{Representação gráfica de \emph{posets}}
  \begin{itemize}
    \item \emph{Posets} podem ser caracterizados de maneira amigável,
        considerando os elementos de $S$ como ``pontos em um plano''
    \item As conexões entre elementos são análogas à descrição das relações, e
        utilizam-se ``vetores'' para simbolizá-las
    \item De maneira formal, um \emph{poset} pode ser representado como um
        grafo dirigido, cujos vértices são os elementos de $S$, e arestas os
          elementos de $\preccurlyeq$
  \end{itemize}
\end{frame}

\begin{frame}
  \frametitle{Diagrama de Hasse}
  \begin{itemize}
    \item Deseja-se desenhar o \emph{poset} com o menor número de conexões
        possível (através da redução transitiva do grafo correspondente)
    \item Esta técnica gera um esquema chamado de \textbf{diagrama de Hasse},
        que representa o \emph{poset} de maneira sucinta
    \item Entretanto, existem várias formas de desenhá-los, tornando-se difícil
        criar diagramas ordenados
    \item O diagrama de Hasse de \emph{tosets} dá origem ao nome de ``cadeia''
  \end{itemize}
\end{frame}

\begin{frame}
  \frametitle{Exemplo}
  \framesubtitle{Diagrama de Hasse}
  \centering
  \resizebox{6cm}{!}{
    \begin{tikzpicture}
      \node [small, label=above:{\tiny $\{x, y, z\}$}] (xyz) at (0, 0) {};
      \node [small, below left=of xyz, label=left:{\tiny $\{x, y\}$}] (xy) {};
      \node [small, below right=of xyz, label=right:{\tiny $\{y, z\}$}] (yz) {};
      \node [small, below=of xyz, label=right:{\tiny $\{x, z\}$}] (xz) {};
      \node [small, below=of xy, label=left:{\tiny $\{x\}$}] (x) {};
      \node [small, below=of xz, label=right:{\tiny $\{y\}$}] (y) {};
      \node [small, below=of yz, label=right:{\tiny $\{z\}$}] (z) {};
      \node [small, below=of y, label=below:{\tiny $\emptyset$}] (0) {};
      \draw[->] (0)   edge                                       (x)
                (0)   edge                                       (y)
                (0)   edge                                       (z)
                (x)   edge                                       (xy)
                (x)   edge                                       (xz)
                (y)   edge                                       (xy)
                (y)   edge                                       (yz)
                (z)   edge                                       (xz)
                (z)   edge                                       (yz)
                (xy)  edge                                       (xyz)
                (xz)  edge                                       (xyz)
                (yz)  edge                                       (xyz)
                (0)   edge                                       (xy)
                (0)   edge [bend right]                          (yz)
                (0)   edge [bend left]                           (xz)
                (0)   edge [bend left]                           (xyz)
                (x)   edge                                       (xyz)
                (y)   edge [bend left]                           (xyz)
                (z)   edge [bend left]                           (xyz)
                (0)   edge [loop, in=225, out=180, looseness=20] (0)
                (x)   edge [loop, in=270, out=225, looseness=20] (x)
                (y)   edge [loop, in=225, out=180, looseness=20] (y)
                (z)   edge [loop, in=270, out=315, looseness=20] (z)
                (xy)  edge [loop, in=90,  out=135, looseness=20] (xy)
                (yz)  edge [loop, in=90,  out=45,  looseness=20] (yz)
                (xz)  edge [loop, in=135, out=90,  looseness=20] (xz)
                (xyz) edge [loop, in=180, out=135, looseness=20] (xyz);
    \end{tikzpicture}
  }
  \begin{itemize}
    \item $(\mathcal{P}(\{x, y, z\}), \subseteq)$ com todas as arestas
        relativas às relações entre elementos
  \end{itemize}
\end{frame}

\begin{frame}
  \frametitle{Exemplo}
  \framesubtitle{Diagrama de Hasse}
  \centering
  \resizebox{6cm}{!}{
    \begin{tikzpicture}
      \node [small, label=above:{\tiny $\{x, y, z\}$}] (xyz) at (0, 0) {};
      \node [small, below left=of xyz, label=left:{\tiny $\{x, y\}$}] (xy) {};
      \node [small, below right=of xyz, label=right:{\tiny $\{y, z\}$}] (yz) {};
      \node [small, below=of xyz, label=right:{\tiny $\{x, z\}$}] (xz) {};
      \node [small, below=of xy, label=left:{\tiny $\{x\}$}] (x) {};
      \node [small, below=of xz, label=right:{\tiny $\{y\}$}] (y) {};
      \node [small, below=of yz, label=right:{\tiny $\{z\}$}] (z) {};
      \node [small, below=of y, label=below:{\tiny $\emptyset$}] (0) {};
      \draw[->] (0)   edge                                       (x)
                (0)   edge                                       (y)
                (0)   edge                                       (z)
                (x)   edge                                       (xy)
                (x)   edge                                       (xz)
                (y)   edge                                       (xy)
                (y)   edge                                       (yz)
                (z)   edge                                       (xz)
                (z)   edge                                       (yz)
                (xy)  edge                                       (xyz)
                (xz)  edge                                       (xyz)
                (yz)  edge                                       (xyz)
                (0)   edge [loop, in=225, out=180, looseness=20] (0)
                (x)   edge [loop, in=270, out=225, looseness=20] (x)
                (y)   edge [loop, in=225, out=180, looseness=20] (y)
                (z)   edge [loop, in=270, out=315, looseness=20] (z)
                (xy)  edge [loop, in=90,  out=135, looseness=20] (xy)
                (yz)  edge [loop, in=90,  out=45,  looseness=20] (yz)
                (xz)  edge [loop, in=135, out=90,  looseness=20] (xz)
                (xyz) edge [loop, in=180, out=135, looseness=20] (xyz);
    \end{tikzpicture}
  }
  \begin{itemize}
    \item As arestas relativas à transitividade da relação são omitidas
  \end{itemize}
\end{frame}

\begin{frame}
  \frametitle{Exemplo}
  \framesubtitle{Diagrama de Hasse}
  \centering
  \resizebox{6cm}{!}{
    \begin{tikzpicture}
      \node [small, label=above:{\tiny $\{x, y, z\}$}] (xyz) at (0, 0) {};
      \node [small, below left=of xyz, label=left:{\tiny $\{x, y\}$}] (xy) {};
      \node [small, below right=of xyz, label=right:{\tiny $\{y, z\}$}] (yz) {};
      \node [small, below=of xyz, label=right:{\tiny $\{x, z\}$}] (xz) {};
      \node [small, below=of xy, label=left:{\tiny $\{x\}$}] (x) {};
      \node [small, below=of xz, label=right:{\tiny $\{y\}$}] (y) {};
      \node [small, below=of yz, label=right:{\tiny $\{z\}$}] (z) {};
      \node [small, below=of y, label=below:{\tiny $\emptyset$}] (0) {};
      \draw[->] (0) edge (x)  (0)  edge (y)   (0)  edge (z)   (x)  edge (xy)
                (x) edge (xz) (y)  edge (xy)  (y)  edge (yz)  (z)  edge (xz)
                (z) edge (yz) (xy) edge (xyz) (xz) edge (xyz) (yz) edge (xyz);
    \end{tikzpicture}
  }
  \begin{itemize}
    \item As arestas relativas à reflexividade da relação também são omitidas
  \end{itemize}
\end{frame}

\begin{frame}
  \frametitle{Exemplo}
  \framesubtitle{Diagrama de Hasse}
  \centering
  \resizebox{6cm}{!}{
    \begin{tikzpicture}
      \node [small, label=above:{\tiny $\{x, y, z\}$}] (xyz) at (0, 0) {};
      \node [small, below left=of xyz, label=left:{\tiny $\{x, y\}$}] (xy) {};
      \node [small, below right=of xyz, label=right:{\tiny $\{y, z\}$}] (yz) {};
      \node [small, below=of xyz, label=right:{\tiny $\{x, z\}$}] (xz) {};
      \node [small, below=of xy, label=left:{\tiny $\{x\}$}] (x) {};
      \node [small, below=of xz, label=right:{\tiny $\{y\}$}] (y) {};
      \node [small, below=of yz, label=right:{\tiny $\{z\}$}] (z) {};
      \node [small, below=of y, label=below:{\tiny $\emptyset$}] (0) {};
      \draw (0) edge (x)  (0)  edge (y)   (0)  edge (z)   (x)  edge (xy)
            (x) edge (xz) (y)  edge (xy)  (y)  edge (yz)  (z)  edge (xz)
            (z) edge (yz) (xy) edge (xyz) (xz) edge (xyz) (yz) edge (xyz);
    \end{tikzpicture}
  }
  \begin{itemize}
    \item O direcionamento é omitido, visto que todas as arestas apontam ``para
        cima''
  \end{itemize}
\end{frame}

\begin{frame}
  \frametitle{Elementos extremos}
  \begin{itemize}
    \item<1-> É razoável conhecer os elementos que denotam as fronteiras de um
        \emph{poset} $(S, \preccurlyeq)$
    \item Considere um elemento qualquer $s' \in S$, e $\forall s \in S$
    \begin{itemize}[itemsep=0pt]
      \item<2-> $s'$ é \textbf{máximo} se é ``maior que'' todos os outros, ou
          seja, $s \prec s'$
      \item<3-> $s'$ é \textbf{mínimo} se é ``menor que'' todos os outros, ou
          seja, $s' \prec s$
      \item<4-> $s'$ é \textbf{maximal} se não existe outro ``maior que'' este,
          ou seja, $s' \preccurlyeq s$
      \item<5-> $s'$ é \textbf{minimal} se não existe outro ``menor que'' este,
          ou seja, $s \preccurlyeq s'$
    \end{itemize}
  \end{itemize}
\end{frame}

\begin{frame}
  \frametitle{Elementos extremos}
  \begin{itemize}
    \item Portanto, devem existir até um elemento máximo e um elemento mínimo
    \item Se este for o caso, estes serão o único elemento maximal e minimal, e
        serão denotados por $\top, \bot$ respectivamente
    \item À ausência destes, não existem restrições no número de elementos
        minimais ou maximais, e estes conjuntos serão denotados por $S_{\min},
          S_{\max}$ respectivamente
    \item Poderão existir elementos maximal e minimal únicos, mas não
        necessariamente serão o máximo e mínimo
  \end{itemize}
\end{frame}

\begin{frame}
  \frametitle{Exemplo}
  \framesubtitle{Elementos extremos}
  \begin{itemize}
    \item<1-> Considere o \emph{poset} $(\mathbb{Z}^{+}, \mid)$. É possível
        identificar elementos extremos?
    \begin{itemize}[itemsep=0pt]
      \item<2-> Como o \emph{poset} tende a $+\infty$, não existem elementos
          maximais e máximo
      \item<3-> $\bot = 1$, pois $1 \mid n, \forall n \in \mathbb{Z}^{+}$
    \end{itemize}
    \item<4-> Considere o \emph{poset} $(\mathcal{P}(T), \subseteq)$. É
        possível identificar elementos extremos?
    \begin{itemize}[itemsep=0pt]
      \item<5-> Tome um subconjunto qualquer $t \in \mathcal{P}(T)$
      \item<6-> Observe que $\emptyset \subseteq t$, e portanto $\bot =
          \emptyset$
      \item<7-> Da mesma maneira, $t \subseteq T$, e portanto $\top = T$
    \end{itemize}
  \end{itemize}
\end{frame}

\begin{frame}
  \frametitle{Exemplo}
  \framesubtitle{Elementos extremos}
  \centering
  \resizebox{6cm}{!}{
    \begin{tikzpicture}
      \node [small, below left=of xyz, label=above:{\tiny $\{x, y\}$}] (xy) {};
      \node [small, below right=of xyz, label=above:{\tiny $\{y, z\}$}] (yz) {};
      \node [small, below=of xyz, label=above:{\tiny $\{x, z\}$}] (xz) {};
      \node [small, below=of xy, label=below:{\tiny $\{x\}$}] (x) {};
      \node [small, below=of xz, label=below:{\tiny $\{y\}$}] (y) {};
      \node [small, below=of yz, label=below:{\tiny $\{z\}$}] (z) {};
      \draw (x) edge (xy) (x) edge (xz) (y) edge (xy)
            (y) edge (yz) (z) edge (xz) (z) edge (yz);
    \end{tikzpicture}
  }
  \begin{itemize}
    \item<1-> Considere o diagrama de Hasse acima. É possível identificar
        elementos extremos?
    \item<2-> $S_{\min} = \{\{x\}, \{y\}, \{z\}\}, S_{\max} = \{\{x, y\}, \{x,
        z\}, \{y, z\}\}$
  \end{itemize}
\end{frame}

\begin{frame}
  \frametitle{Subconjuntos de \emph{posets}}
  \begin{itemize}
    \item Em um \emph{poset} qualquer $(S, \preccurlyeq)$, a relação é
        preservada para todos os subconjuntos de $S$
    \item Elementos extremos para estes subconjuntos têm nomes especiais
    \begin{itemize}[itemsep=0pt]
      \item Tome um subconjunto $K \subseteq S$ e elemento $x \in S$ quaisquer
      \item $x$ está na \textbf{cota superior} de $K$ se $k \preccurlyeq x,
          \forall k \in K$, e na \textbf{cota inferior} de $K$ se $x
            \preccurlyeq k, \forall k \in K$
      \item Não existe notação definida para as cotas de $K$ em relação à $S$
    \end{itemize}
  \end{itemize}
\end{frame}

\begin{frame}
  \frametitle{Elementos extremos em cotas}
  \begin{itemize}
    \item Dentro da cota superior para um subconjunto de um \emph{poset},
        pode existir um elemento mínimo neste, e este é chamado de
          \textbf{menor cota superior}, ou ínfimo --- $\inf(K)$
    \item Analogamente, pode existir um elemento máximo na cota inferior,
        chamado de \textbf{maior cota inferior}, ou supremo --- $\sup(K)$
    \item Se $\inf(K)$ existe e pertence a $K$, então é um elemento
        minimal ou mínimo de $K$, e de maneira dual, $\sup(K)$ é um
          elemento maximal ou máximo de $K$
  \end{itemize}
\end{frame}

\begin{frame}
  \frametitle{Exemplo}
  \framesubtitle{Elementos extremos em cotas}
  \centering
  \begin{tikzpicture}
    \node [small, label=left:{\tiny  6}]  (6) at (-0.5, -1.6) {};
    \node [small, label=right:{\tiny 10}] (10) at (0.5, -1.6) {};
    \node [small,                      label=above:{\tiny 60}] (60) {};
    \node [small, below left  = of 60, label=left:{\tiny 12}] (12) {};
    \node [small, below       = of 60, label=right:{\tiny 20}] (20) {};
    \node [small, below right = of 60, label=right:{\tiny 30}] (30) {};
    \node [small,       left  = of  6, label= left:{\tiny  4}] (4) {};
    \node [small,       right = of 10, label=right:{\tiny 15}] (15) {};
    \node [small,       below = of 12, label=left:{\tiny  2}] (2) {};
    \node [small,       below = of 20, label=left:{\tiny  3}] (3) {};
    \node [small,       below = of 30, label=right:{\tiny  5}] (5) {};
    \node [small,       below = of  3, label=below:{\tiny  1}] (1) {};
    \draw  (2) edge  (4)  (2) edge  (6)  (2) edge (10)  (3) edge  (6)
           (3) edge (15)  (4) edge (12)  (4) edge (20)  (5) edge (10)
           (5) edge (15)  (6) edge (12)  (6) edge (30) (10) edge (20)
          (10) edge (30) (15) edge (30) (12) edge (60) (20) edge (60)
          (30) edge (60)  (1) edge  (2)  (1) edge  (3)  (1) edge  (5);
  \end{tikzpicture}
  \begin{itemize}
    \item<1-> Considere o \emph{poset} de divisores de $60$ e $K = \{6, 15\}$
    \begin{itemize}[itemsep=0pt]
      \item<2-> A cota superior é $\{30, 60\}$, a cota inferior é $\{1, 3\}$,
          $\inf(K) = 30$ e $\sup(K) = 3$
    \end{itemize}
  \end{itemize}
\end{frame}

\begin{frame}
  \frametitle{Exemplo}
  \framesubtitle{Elementos extremos em cotas}
  \centering
  \begin{tikzpicture}
      \node [small,                    label=above:{\tiny h}] (h) {};
      \node [small, right      = of h, label=above:{\tiny j}] (j) {};
      \node [small, below      = of j, label=right:{\tiny f}] (f) {};
      \node [small, below      = of f, label=right:{\tiny e}] (e) {};
      \node [small, below      = of e, label=right:{\tiny c}] (c) {};
      \node [small, below left = of h, label=left:{\tiny g}] (g) {};
      \node [small, below      = of g, label=left:{\tiny d}] (d) {};
      \node [small, below      = of d, label=left:{\tiny b}] (b) {};
      \node [small, below left = of c, label=below:{\tiny a}] (a) {};
      \draw (a) edge (b) (a) edge (c) (b) edge (d) (d) edge (g) (g) edge (h)
            (c) edge (e) (e) edge (f) (f) edge (j) (f) edge (h) (d) edge (f)
            (b) edge (e);
  \end{tikzpicture}
  \begin{itemize}
    \item<1-> Considere o \emph{poset} acima e $K = \{a, c, d, f\}$
    \begin{itemize}[itemsep=0pt]
      \item<2-> A cota superior é $\{f, h, j\}$, a cota inferior é $\{a\}$,
          $\inf(K) = f$ e $\sup(K) = a$
    \end{itemize}
  \end{itemize}
\end{frame}

\begin{frame}
  \frametitle{Material de estudo}
  \nocite{*}
  \bibliographystyle{plain}
  \bibliography{ref}
  \begin{itemize}
    \item Leitura das páginas 618-626 e resolução dos exercícios 1-2, 5-6,
        9-11, 13-16, 18, 20-25, 32-35, nas páginas 630-631.
  \end{itemize}
\end{frame}

\end{document}
